\section{Introduction}
\label{sec:intro}

No quantum system is truly isolated.  A qubit in a device interacts with control lines, material defects, and electromagnetic modes; an electronic excitation in a molecule interacts with nuclear motion and its solvent; and an emitter in a solid exchanges energy with photons and lattice vibrations.  Open quantum-system theory begins by drawing a boundary around the degrees of freedom of interest---the \emph{system}---and calling everything outside that boundary the \emph{environment} or \emph{bath}.  If the joint state is $\rho_{SE}(t)$, the state available to an observer of the system is
\begin{equation}
    \rho_S(t)=\operatorname{Tr}_E\!\left[\rho_{SE}(t)\right].
\end{equation}
This reduction is useful precisely because the environment is normally too large to simulate or observe in full.  It is also incomplete: the partial trace discards which environmental degrees of freedom became correlated with the system and whether those correlations can influence it again.

In a Markovian approximation, the omitted information is assumed to become irrelevant on the timescale resolved by the system.  The environment may then be represented by a time-local generator or by a sequence of memoryless channels.  Such an approximation is well motivated in regimes with weak coupling and a clear separation between environmental and system timescales, but an exponentially decaying correlation function by itself is not a universal guarantee of Markovian dynamics~\cite{Rivas2014,deVega2017}.  Memory must instead be retained when the environment responds on comparable timescales, when the coupling is strong, when the reservoir has narrow spectral structure or finite extent, or when emitted excitations can return to the system.  Examples include time-delayed feedback in photonic waveguides~\cite{Pichler2016}, femtosecond exciton--phonon dynamics in semiconductor quantum dots~\cite{Liu2019}, and electronic energy transfer coupled to long-lived molecular vibrations~\cite{Chin2013}.  In such settings, discarding environmental memory can change qualitative as well as quantitative predictions.  Recent reviews describe the breadth of these regimes and the non-perturbative tensor-network methods developed for them~\cite{Keeling2026,LacroixReview2026}.

Quantum mechanics makes the missing information especially consequential.  Suppose that two preparation histories lead to the same reduced state $\rho_S(t)$ but to different joint states $\rho_{SE}(t)$.  The system looks identical at time $t$, yet its correlations with the environment differ, so the same subsequent control operation can lead to different future statistics.  Moreover, an intervention on the system---a preparation, unitary pulse, noisy channel, or measurement---acts on the joint state and can condition the environment on its outcome.  A trajectory of reduced states, or a two-time dynamical map inferred for one preparation procedure, therefore need not predict what happens under a different sequence of interventions.  This is not a consequence of merely \emph{knowing} the state: it is the physical action of the intervention on a correlated system--environment state that matters.  Experimental process characterization on quantum processors has demonstrated the operational relevance of these multi-time correlations~\cite{White2020}.

The process tensor supplies the required multi-time description~\cite{Pollock2018Framework,Pollock2018Markov,Milz2021}.  For intervention times $t_0<\cdots<t_{n-1}$, let $\mathcal{A}_{x_k}^{(k)}$ denote the completely positive map associated with outcome $x_k$ of an instrument applied at $t_k$.  A process tensor $\mathcal{T}_{n:0}$ maps the ordered intervention sequence to an unnormalised final state,
\begin{equation}
    \widetilde{\rho}_n(\boldsymbol{x})
    =\mathcal{T}_{n:0}
      \!\left[
      \mathcal{A}_{x_{n-1}}^{(n-1)},\ldots,
      \mathcal{A}_{x_0}^{(0)}
      \right],
    \qquad
    p(\boldsymbol{x})=\operatorname{Tr}\widetilde{\rho}_n(\boldsymbol{x}).
    \label{eq:process-tensor-action}
\end{equation}
It is multilinear in the interventions and, in its Choi representation, is a positive operator whose input and output spaces are ordered in time and constrained by causality.  This representation connects process tensors to quantum combs and general quantum networks~\cite{Chiribella2009}.  Markovian dynamics is contained as a special case in which the multi-time process factorises into independent channels.  Conversely, a general process tensor retains precisely the information required to predict the result of arbitrary admissible interventions, including temporally correlated experiments described by testers with persistent ancillary systems.

The formal completeness of a process tensor creates its main numerical difficulty.  A dense Choi representation acquires an input--output pair of system indices at every time step and therefore grows exponentially with the number of interventions.  Representing the same object as a process-tensor matrix product operator (PT-MPO) reorganises this exponential tensor into a one-dimensional network along time.  Its physical legs accept the interventions, while its virtual bonds carry correlations across temporal cuts.  Singular-value truncation can then compress weakly occupied temporal correlation channels.  Once constructed, the PT-MPO can be reused for many initial states, controls, instruments, and observables: the environmental influence is built once and interrogated many times.

This compression is powerful but not automatic.  The bond dimension may grow rapidly when temporal correlations are strong or complex, and no PT-MPO representation is guaranteed to remain efficient for arbitrary dynamics.  Discretisation, bath truncation, and singular-value thresholds also introduce distinct convergence requirements.  Thus, a process tensor is a unifying physical object, not a promise that every environment is computationally easy; practical performance depends on the bath, the time window, the target observables, and the algorithm used to construct and contract the tensor~\cite{Keeling2026,LacroixReview2026}.

Several tensor-network constructions address different parts of this problem.  The time-evolving matrix-product-operator (TEMPO) algorithm compresses a discretised Feynman--Vernon influence functional for Gaussian environments~\cite{FeynmanVernon1963,Strathearn2018}.  PT-TEMPO contracts this influence network into a reusable causal PT-MPO~\cite{Jorgensen2019}.  Automated compression of environments (ACE) instead constructs the influence of a collection of independent environmental modes and combines those influences while repeatedly compressing the temporal bonds; because the individual modes may be finite-dimensional and non-Gaussian, ACE covers models beyond the Gaussian-bosonic setting~\cite{Cygorek2022}.  A complementary family of methods, including TEDOPA, maps a continuous environment to a chain and evolves an enlarged system--environment state as a matrix product state~\cite{Chin2010,Prior2010}.  These formulations can be understood as different representations or contraction strategies for the underlying system--environment influence, and none is uniformly optimal.

The corresponding software ecosystem is already substantial.  OQuPy provides Python implementations of PT-TEMPO and process-tensor workflows for Gaussian bosonic environments, together with multi-time correlations, control, multiple environments, and process reuse~\cite{Fux2024}.  The C++ ACE toolkit implements ACE and a wider collection of PT-MPO construction and composition methods, including PT-TEMPO and divide-and-conquer schemes~\cite{CygorekGauger2024}.  QuantumDynamics.jl offers a modular Julia interface to several open-system techniques, including real-time path-integral and influence-functional methods~\cite{Bose2023}.  MPSDynamics.jl follows the explicit-environment route in Julia, implementing finite-temperature T-TEDOPA with matrix-product and tree-tensor-network states~\cite{Lacroix2024}.  These projects are complementary rather than interchangeable: they differ in bath assumptions, numerical representation, supported observables, and programming abstractions.

\texttt{ProcessTensors.jl}\footnote{\url{https://github.com/Gauthameshwar/ProcessTensors.jl}} occupies a distinct role in this landscape.  It is an ITensor- and ITensorMPS-native framework in which the PT-MPO is a first-class reusable object, rather than an internal by-product tied to one solver.  ITensor's index-aware interface permits code to follow tensor diagrams through named and tagged legs instead of depending on a fragile storage order~\cite{Fishman2022}.  \texttt{ProcessTensors.jl} extends this approach to Liouville space and separates three tasks: constructing an environmental process, defining the experiment contracted with it, and characterising the temporal structure that remains.  The scientific contribution is therefore an interoperable software architecture and operational toolkit, not a new principle for constructing process tensors.

The present release realises this goal through five connected capabilities.  An exact dense builder treats small system--environment models and provides a transparent reference against which compressed methods can be tested.  A scalable ACE builder produces compressed PT-MPOs for environments assembled from independent finite-dimensional modes, with explicit control of the time step, singular-value cutoff, and maximum bond dimension.  Both builders return the same process-tensor abstraction.  An instrument gallery supplies common preparations, controls, channels, and measurements as composable operations, while a tester interface supports experiments whose ancillary memory persists between intervention times.  Finally, temporal-bond routines expose the Schmidt spectrum and its operator-space entanglement entropy (OSEE) across a selected PT-MPO cut~\cite{Prosen2007}.  We use OSEE as a representation-dependent diagnostic of temporal correlation structure and compressibility, not as a universal scalar measure of non-Markovianity.

The package is intended to be useful at three levels.  For researchers already working with ITensor, it provides a direct route from familiar MPS/MPO operations to non-Markovian multi-time calculations.  For researchers entering the process-tensor field, the dense builder, instrument gallery, and documented examples provide a path from the defining equations to executable contractions.  For method developers, a common builder interface, tagged-index conventions, and exact small-model references provide infrastructure on which further influence-functional, chain-mapping, periodic, infinite, or tree-geometry constructions can be developed.  The emphasis is on making process tensors inspectable and composable while preserving access to the underlying ITensor objects.

Validation and benchmarking are consequently part of the package's scientific claim.  We test index and Liouville-space conventions, causality and trace-preservation conditions, agreement with exact diagonalisation, and agreement between dense and compressed builders where both are feasible.  Benchmarks report runtime, peak memory, and maximum temporal bond dimension alongside errors in final states and selected multi-time observables.  Comparisons with the C++ ACE toolkit and PT-TEMPO are made on matched Hamiltonians, discretisations, bath truncations, and convergence targets.  Their purpose is to establish correctness and expose practical trade-offs, rather than to claim universal superiority of one algorithm or implementation.

The remainder of this paper is organised as follows.  We first introduce the process-tensor, Choi, and Liouville-space conventions and connect their tensor diagrams to tagged ITensor indices.  We then describe the package architecture and the dense and ACE construction backends.  The following sections develop instruments, testers, and temporal-bond diagnostics through reproducible examples.  We next present correctness tests and performance benchmarks against independent implementations.  We conclude by delineating the present domain of validity and discussing additional PT-MPO builders, dissipative process tensors, and richer temporal-correlation diagnostics as future directions.
